% !TEX program = pdflatex
% Raport technologiczny projektu „Sąsiedzko" (HackYeah 2026, kategoria Smart City).
% Samodzielny plik — wklej w całości do Overleaf (silnik: pdfLaTeX).
\documentclass[11pt,a4paper]{article}

\usepackage[T1]{fontenc}
\usepackage[utf8]{inputenc}
\usepackage[polish]{babel}
\usepackage{lmodern}
\usepackage[a4paper,margin=2.5cm]{geometry}
\usepackage{booktabs}
\usepackage{tabularx}
\usepackage{enumitem}
\usepackage{xcolor}
\usepackage[hidelinks]{hyperref}
\usepackage{titlesec}

\definecolor{akcent}{HTML}{6D28D9}
\titleformat{\section}{\Large\bfseries\color{akcent}}{\thesection.}{0.5em}{}
\titleformat{\subsection}{\large\bfseries}{\thesubsection.}{0.5em}{}
\setlist{itemsep=2pt,topsep=3pt}

\title{\textbf{Sąsiedzko} \\[0.3em] \large Raport technologiczny}
\author{Zespół HackYeah 2026 --- kategoria otwarta Smart City}
\date{Październik 2026}

\begin{document}
\maketitle

\section{Wprowadzenie}

\textbf{Sąsiedzko} to instalowalna aplikacja webowa (PWA), w której mieszkańcy dodają
sąsiedzkie wydarzenia na mapie miasta, zapisują się na nie jednym kliknięciem i rozmawiają
na czatach w czasie rzeczywistym. Niniejszy raport opisuje technologie użyte do budowy
projektu oraz --- zgodnie z wymogami regulaminu --- uzasadnienie każdej decyzji.

Przyjęliśmy zasadę \emph{prostego, możliwego do obrony stosu}: minimum warstw, przewidywalne
narzędzia i jak najmniej własnej infrastruktury, aby w czasie hackathonu skupić się na funkcjach,
a nie na konfiguracji serwerów.

\section{Przegląd stosu technologicznego}

\begin{table}[h]
\centering
\renewcommand{\arraystretch}{1.3}
\begin{tabularx}{\textwidth}{@{}l X@{}}
\toprule
\textbf{Warstwa} & \textbf{Technologia} \\
\midrule
Frontend        & React 19 + Vite 8 + TypeScript + MUI (Material UI) 9 \\
Mapa            & Leaflet + OpenStreetMap, grupowanie pinezek (react-leaflet-cluster) \\
Backend / dane  & Supabase --- PostgreSQL, Auth (magic link), Realtime, Storage \\
AI              & Supabase Edge Function (Deno) $\rightarrow$ Claude (Anthropic) \\
Automatyzacja   & Bot w Node.js uruchamiany przez GitHub Actions \\
Hosting         & Render (Static Site) \\
PWA             & Manifest + ikony (instalacja na telefonie) \\
Jakość kodu     & TypeScript (tsc), oxlint \\
\bottomrule
\end{tabularx}
\caption{Warstwy aplikacji i użyte technologie.}
\end{table}

\section{Frontend}

\begin{itemize}
  \item \textbf{React 19} --- biblioteka do budowy interfejsu w architekturze komponentowej.
        Wybrana ze względu na dojrzałość, ogromny ekosystem i znajomość w zespole.
  \item \textbf{Vite 8} --- bundler i serwer deweloperski. Błyskawiczny start, natychmiastowy
        hot reload i prosta konfiguracja --- kluczowe przy pracy pod presją czasu.
  \item \textbf{TypeScript} --- statyczne typowanie. Wspólne typy danych (\texttt{src/lib/types.ts})
        odpowiadają tabelom w bazie, co ogranicza błędy na styku frontend--backend.
  \item \textbf{MUI (Material UI) 9} --- gotowe, dostępne komponenty z systemem motywów.
        Pozwolił szybko zbudować spójny, mobilny interfejs zgodny z WCAG.
  \item \textbf{Emotion} --- silnik stylów (CSS-in-JS) używany przez MUI.
  \item \textbf{React Router 7} --- routing po stronie klienta (mapa, dodawanie, czat, konto, panel administratora).
  \item \textbf{Web Speech API} --- wbudowane w przeglądarkę rozpoznawanie mowy; służy do
        dyktowania zdania w funkcji „wydarzenie z jednego zdania" (bez zewnętrznej biblioteki).
  \item \textbf{Wielojęzyczność (PL/EN)} i animacje --- rozwiązania własne, oparte na lekkim
        kontekście React i CSS, bez dodatkowych zależności.
\end{itemize}

\section{Mapa}

\begin{itemize}
  \item \textbf{Leaflet} wraz z \textbf{react-leaflet} --- lekka, otwarta biblioteka map.
        W przeciwieństwie do rozwiązań komercyjnych nie wymaga klucza API ani opłat.
  \item \textbf{OpenStreetMap} --- otwarte źródło kafelków mapy (dane tworzone społecznościowo).
  \item \textbf{react-leaflet-cluster} --- grupowanie pinezek: wiele bliskich wydarzeń łączy się
        w jeden znacznik z licznikiem i rozwija po przybliżeniu, co utrzymuje czytelność mapy.
\end{itemize}

\section{Backend i dane --- Supabase}

Zamiast pisać własny serwer wykorzystaliśmy \textbf{Supabase} --- platformę typu
\emph{Backend-as-a-Service} opartą na PostgreSQL.

\begin{itemize}
  \item \textbf{PostgreSQL} --- relacyjna baza danych; schemat i migracje trzymamy w repozytorium
        (\texttt{supabase/schema.sql}, \texttt{supabase/migrations/}).
  \item \textbf{Row Level Security (RLS)} --- bezpieczeństwo danych egzekwowane na poziomie bazy.
        Np. dokładny adres wydarzenia jest widoczny dopiero po zapisie, a spotkania „we dwoje"
        baza zwraca wyłącznie zalogowanym osobom pełnoletnim. Reguł nie da się obejść z frontu.
  \item \textbf{Auth (magic link)} --- logowanie bez haseł (imię + e-mail). Mniej danych do ochrony
        i niższa bariera wejścia dla użytkownika.
  \item \textbf{Realtime} --- subskrypcja zmian w tabeli \texttt{messages} zasila czat na żywo
        bez odpytywania serwera.
  \item \textbf{Storage} --- przechowywanie zdjęć wydarzeń (do 5 MB na plik).
\end{itemize}

\section{Sztuczna inteligencja (AI)}

\subsection{Wydarzenie z jednego zdania}
Główny wyróżnik produktu. Użytkownik pisze lub dyktuje jedno zdanie (np. \emph{„urodziny babci
w sobotę o 15 w parku Jordana, do 10 osób, dla seniorów"}), a formularz wypełnia się sam.

\begin{itemize}
  \item \textbf{Claude (Anthropic)} --- model językowy (LLM) zamieniający zdanie na pola wydarzenia.
  \item \textbf{Supabase Edge Function (Deno)} --- funkcja brzegowa \texttt{parse-event} pośredniczy
        między aplikacją a modelem. Dzięki temu \textbf{klucz API nigdy nie trafia do przeglądarki}
        --- trzyma się wyłącznie w sekretach funkcji.
  \item \textbf{Tool use ze ścisłym schematem} --- model wywołuje narzędzie o zdefiniowanym schemacie,
        w którym kategoria i grupy docelowe są ograniczone do listy z bazy (\texttt{enum}). Dzięki temu
        model nie wymyśla wartości spoza dozwolonego zbioru, a wynik jest przewidywalny.
\end{itemize}

\subsection{Świadome granice użycia AI}
Klasyfikację wydarzeń importowanych przez bota (kategoria + grupy docelowe) wykonują
\textbf{reguły słów kluczowych}, a nie model AI --- celowo, ponieważ są przewidywalne, darmowe
i nie wymagają wywołań sieciowych. AI stosujemy tam, gdzie realnie podnosi jakość (język naturalny),
a nie wszędzie „na siłę".

\subsection{AI we wsparciu programisty}
Do pisania kodu korzystaliśmy z asystenta \textbf{Claude Code} (Anthropic). Każdą decyzję
architektoniczną rozumiemy i potrafimy obronić.

\section{Automatyzacja --- bot importujący wydarzenia}

\begin{itemize}
  \item \textbf{Node.js} --- skrypt bota (\texttt{npm run bot}) raz dziennie losuje kilka wydarzeń
        z miejskiego kalendarza \textbf{Karnet} (karnet.krakowculture.pl), dobiera im kategorię
        i grupy docelowe regułami, a następnie dodaje do bazy bez duplikatów i z linkiem do oryginału.
  \item \textbf{GitHub Actions} --- harmonogram (cron) uruchamiający bota; sekrety przechowywane
        w ustawieniach repozytorium, nigdy w kodzie.
\end{itemize}

\section{PWA, hosting i jakość kodu}

\begin{itemize}
  \item \textbf{PWA} --- manifest i zestaw ikon umożliwiają instalację aplikacji na ekranie głównym
        telefonu (bez trybu offline).
  \item \textbf{Render (Static Site)} --- hosting zbudowanej aplikacji. Krok budowania kopiuje
        \texttt{index.html} do folderów podstron, aby odświeżenie adresu typu \texttt{/mapa/}
        działało bez dodatkowych reguł serwera.
  \item \textbf{Dostępność (WCAG 2.1)} --- domyślny wygląd spełnia poziom AA (kontrast, widoczny fokus,
        obsługa klawiaturą), a tryb „Duży i wyraźny widok" podnosi kontrast do AAA i powiększa elementy.
  \item \textbf{tsc} (sprawdzanie typów) i \textbf{oxlint} (szybki linter) --- utrzymanie jakości kodu.
\end{itemize}

\section{Źródła danych}

\begin{itemize}
  \item \textbf{OpenStreetMap} --- kafelki mapy.
  \item \textbf{Karnet} (Krakowskie Biuro Festiwalowe) --- publiczna lista wydarzeń importowanych przez bota.
  \item \textbf{Dane demonstracyjne} --- wydarzenia wymyślone na potrzeby prezentacji, osadzone
        w prawdziwych lokalizacjach; brak prawdziwych danych osobowych.
\end{itemize}

\section{Podsumowanie}

Stos projektu „Sąsiedzko" dobraliśmy pod trzy cele: \textbf{szybkość wytwarzania} (React + Vite + MUI,
Supabase zamiast własnego backendu), \textbf{bezpieczeństwo danych} (RLS w bazie, klucz AI w sekretach
funkcji brzegowej) oraz \textbf{brak kosztów i zależności komercyjnych} (OpenStreetMap, Leaflet,
logowanie bez haseł). AI stosujemy punktowo --- tam, gdzie przetwarzanie języka naturalnego realnie
skraca drogę użytkownika do dodania wydarzenia.

\end{document}
